% ChargeShield-FL --- DSN 2027 --- 6. Discussion and Limitations
% Scritta il 2026-09-23. La tabella delle conclusioni consentite traduce la
% guida sez. 11 ("Espressioni da non usare senza l'evidenza corrispondente").

\section{Discussion and Limitations}
\label{sec:discussion}

\subsection{What the Evidence Allows}
\label{sec:disc-allowed}

Table~\ref{tab:allowed} pairs each observation of Section~\ref{sec:results}
with the conclusion it supports and with the stronger reading it does not.

\begin{table*}[t]
\centering
\caption{Conclusions supported by the current evidence, and readings they do
not support.}
\label{tab:allowed}
\footnotesize
\begin{tabular}{@{}p{4.3cm}p{6.1cm}p{5.6cm}@{}}
\toprule
Observation & Supported conclusion & Not supported \\
\midrule
Aggregate AUC near $0.5$ without DP & no attack we ran separates members from non-members on average, within the equivalence margin & the system does not leak membership \\
Per-record excess without DP ($z = 7.56$) & about 120 sessions are ranked systematically high by reconstruction error across seeds & an operational attacker identifies these sessions \\
No excess at $\varepsilon \le 1$, loss $>100\times$ & at these settings the model no longer serves its purpose; protection and loss of utility are confounded & DP eliminated the leakage \\
Canary control passes on Office~1 & the instrument detects induced memorisation in that regime & the instrument validates the natural-regime null \\
Incidents of Table~\ref{tab:incidents} & charging records are personal data and central stores fail in practice & FL or DP would have prevented these incidents \\
\bottomrule
\end{tabular}
\end{table*}

Two points deserve emphasis. The aggregate null and the per-record excess are
not in contradiction: the AUC averages over thousands of sessions, and an
excess of about 120 sessions among 17{,}561 barely moves it. This is exactly
the situation in which reading the AUC alone would report ``no leakage'' while
some records are distinguishable. Conversely, the disappearance of the excess
under DP at $\varepsilon \le 1$ is the result that is easiest to overstate, and
it is the one the $\varepsilon$ sweep is designed to resolve.

\subsection{Relation to the Incidents}
\label{sec:disc-incidents}

The incidents of Section~\ref{sec:threat-incidents} motivate the problem but
are neither evidence for nor against our findings. DP bounds what a trained
model or an update reveals about a protected unit; it does not protect the
databases, cloud platforms and payment interfaces through which the reported
data were lost. FL reduces the need for a central training pool, but each
operator still holds its own records, which were the object of the 2023 and
2024 exposures. The question our measurements answer is therefore
complementary to incident response: if an operator joins a federation, how
much do the artifacts it sends add to what it already has to protect.

\subsection{Limitations}
\label{sec:disc-limits}

\paragraph{Attack calibration and access} The calibrated LiRA degenerates in
this regime, so the per-record analysis is a ranking by loss, and its shadow
models use data that an aggregator would not hold. Our measurements are
upper bounds for the attacks considered, not a demonstration of what a
realistic adversary achieves, and an adaptive attack tailored to this system
is not evaluated.

\paragraph{Meaning of $\varepsilon$} The client-level $\varepsilon$ is a
calibration parameter: the classical Gaussian calibration does not certify
$\varepsilon \ge 1$, normalisation statistics and batch-normalisation buffers
are released outside the mechanism, and the record-level accountant has a
known error. We compare the two protected units at comparable cost, never at
equal $\varepsilon$.

\paragraph{Unit and person} Neither protected unit bounds the contribution of a
person with many sessions or with sessions at several sites, and the
identifiers that let us measure this cover 34.8\%--93.5\% of the sessions
depending on the site.

\paragraph{Scope of the control} The canary control covers one site, in a
regime altered on purpose (duplicated templates, higher capacity, a nearly
unique feature, 1{,}000 epochs). It shows sensitivity there, not at the effect
sizes of the natural regime.

\paragraph{Statistical power} Five seeds per cell bound what paired tests can
show; we report per-seed values and use equivalence tests only with a stated
margin, whose justification is still open.

\paragraph{Utility} We measure reconstruction error on natural data, not the
detector's ability to detect attacks, for which no labels exist. A small loss
does not show that the detector is useful.

\paragraph{Transferability} Three real sites are not necessarily three
independent commercial operators, the model is one small autoencoder, and one
dataset is used. An adapter for a second charging dataset exists, but a single
test run does not support any claim of transfer.

\paragraph{Missing arms} The clip-only arm (B1) does not exist yet, so the
effect of client-level DP cannot be split between clipping and noise; RQ3 and
RQ2 are \pending{pending}.
